\subsection{手动输入公式}

\begin{table}[ht]
    \centering
    \begin{tabular}{|c|c|c|}
    \hline
    运算符类型 & 算符 & 说明 \\ \hline
    数学运算符 & $+$、$-$、$*$、$/$、\^ \ \% & 加、减、乘、除、乘幂、百分比 \\ \hline
    负号 & - & 负号 \\ \hline
    引用运算符 & 冒号:、单个空格、逗号, & 引用公式范围、分隔作用 \\ \hline
    文本运算符 & \& & 连接两个文本字符串 \\ \hline
    比较运算符 & =、<>、<、<=、>、>= & \makecell{等于、不等于、小于、\\ 小于等于、大于、大于等于} \\ \hline
    \end{tabular}
    \end{table}

    \begin{enumerate}
        \item 引用单元格计算
        \item 直接计算
    \end{enumerate}

\subsection{使用公式函数}

函数名(函数参数，函数参数，$\cdots$ ，函数参数) = 返回值

聚合函数：
\begin{itemize}
    \item SUM 求和
    \item AVERAGE 求平均值
    \item COUNT (只对数字单元格计数) COUNTA(数字文本全计数)
    \item MAX 求最大值
    \item MIN 求最小值
\end{itemize}

\subsection{引用}

\begin{itemize}
    \item 相对引用：对于A列第一行的单元格，相对引用是A1。
    \item 绝对引用：列号和行号前面都加了一个美元的符号\$A\$1。
    \item 混合引用：在行号或者列号其中一个前面加了美元符号，A\$1 或 \$A1，表示只固定行或者列。
\end{itemize}

\newpage
\section{文本函数}

\subsection{EXACT函数}

EXACT函数比较两个文本字符串，如果它们完全相同，则返回 TRUE，否则返回 FALSE。 函数 EXACT 区分大小写，但忽略格式上的差异。 使用 EXACT 可以检验在文档中输入的文本。

语法：EXACT(text1, text2)

EXACT 函数语法具有下列参数：
\begin{enumerate}
  \item text1    必需。 第一个文本字符串。
  \item text2    必需。 第二个文本字符串。
\end{enumerate}

\subsection{LEFT、LEFTB函数}

LEFT 从文本字符串的第一个字符开始返回指定个数的字符。

LEFTB 基于所指定的字节数返回文本字符串中的第一个或前几个字符。

\begin{itemize}
  \item Left函数把全角（如汉字）和半角字符（如数字与字母）都计作一个字符
  \item LeftB函数把全角字符计作两个字节、半角字符计作一个字节
\end{itemize}

语法：

\verb|LEFT(text, [num_chars])|

\verb|LEFTB(text, [num_bytes])|


\subsection{RIGHT、RIGHTB函数}

RIGHT 根据所指定的字符数返回文本字符串中最后一个或多个字符。

RIGHTB 根据所指定的字节数返回文本字符串中最后一个或多个字符。

\subsection{MID、MIDB函数}

MID 返回文本字符串中从指定位置开始的特定数目的字符，该数目由用户指定。

MIDB 根据您指定的字节数，返回文本字符串中从指定位置开始的特定数目的字符。

语法：

\verb|MID(text, start_num, num_chars)|

\verb|MIDB(text, start_num, num_bytes)|

\subsection{LEN、LENB 函数}

LEN 返回文本字符串中的字符个数。

LENB 返回文本字符串中用于代表字符的字节数。

语法：

LEN(text)

LENB(text)

\subsection{TEXTJOIN 函数}

TEXTJOIN 函数将多个区域和/或字符串的文本组合起来，并包括你在要组合的各文本值之间指定的分隔符。 如果分隔符是空的文本字符串，则此函数将有效连接这些区域。

语法：
\verb|TEXTJOIN(delimiter, ignore_empty, text1, [text2], …)|

\begin{itemize}
  \item delimiter：分隔符(必需)，文本字符串，或者为空，或用双引号引起来的一个或多个字符，或对有效文本字符串的引用。 如果提供一个数字，则将被视为文本。
  \item \verb|ignore_empty| (必需)，如果为 TRUE，则忽略空白单元格。
  \item \verb|text1|(必需)，要联接的文本项。 文本字符串或字符串数组，如单元格区域中。
  \item \verb|[text2， ...]|(可选)，要联接的其他文本项。 文本项最多可以包含 252 个文本参数 text1。 每个参数可以是一个文本字符串或字符串数组，如单元格区域。
\end{itemize}

例如，\verb|=TEXTJOIN(" ",TRUE,"The","sun","will","come","up","tomorrow")| 将返回The sun will come up tomorrow

\subsection{TEXTSPLIT 函数}

使用列和行分隔符拆分文本字符串。

TEXTSPLIT 函数的工作方式与文本转列向导相同，但采用公式形式。 它允许跨列拆分或按行向下拆分。 它是 TEXTJOIN 函数的反函数。 

语法：

\verb|=TEXTSPLIT(text,col_delimiter,[row_delimiter],[ignore_empty], |

\verb|[match_mode], [pad_with])|

TEXTSPLIT 函数语法具有下列参数：
\begin{itemize}
  \item text       要拆分的文本。 必需。
  \item \verb|col_delimiter|       标记跨列溢出文本的点的文本。
  \item \verb|row_delimiter|       标记向下溢出文本行的点的文本。 可选。
  \item \verb|ignore_empty|       指定 TRUE 以忽略连续分隔符。 默认为 FALSE，将创建一个空单元格。 可选。
  \item \verb|match_mode|    指定 1 以执行不区分大小写的匹配。 默认为 0，这会执行区分大小写的匹配。 可选。
  \item \verb|pad_with|           用于填充结果的值。 默认值为 \verb|#N/A|。
\end{itemize}